Let \(\phi'=T\phi\) and \(\theta'=T^{-\top}\theta\), where \(T\in\operatorname{GL}_d(\mathbb R)\). Then
\[
\phi'(x,a)^\top\theta'=\phi(x,a)^\top T^\top T^{-\top}\theta=\phi(x,a)^\top\theta.
\]
Thus the set of prediction functions representable in the primed coordinates is exactly the same subset \(\mathcal F\subseteq[0,1]^S\). The empirical loss \(L_n\) depends only on prediction values and the data, so its minimizer set is unchanged. The sequential lexicographic rule orders the values of functions on the same fixed ordered set \(S\), rather than coefficient representatives, and therefore the selected function \(\widehat f\) is unchanged.

Likewise, \(\|h\|_n^2=n^{-1}\sum_i h(X_i,A_i)^2\) is coordinate free. Hence \(\mathcal Q_n\) is the same compact subset of prediction space after reparameterization. For every \((x,a)\in S\), the scalar optimization problem
\[
\underline r_n(x,a)=\min_{f\in\mathcal Q_n}f(x,a)
\]
has the same feasible prediction functions and the same objective values in the two coordinate systems. Its minimum value is therefore unchanged. Consequently every numerator and denominator in
\[
\widehat\pi(a\mid x)=\frac{\pi_{\mathrm{ref}}(a\mid x)e^{\eta\underline r_n(x,a)}}{\sum_b\pi_{\mathrm{ref}}(b\mid x)e^{\eta\underline r_n(x,b)}}
\]
is unchanged, so the deployed policy is unchanged at every context, including contexts absent from the sample. Finally, \(\underline J_n\) is written entirely in terms of prediction values in \(\mathcal Q_n\), the observed contexts, \(\pi_{\mathrm{ref}}\), and \(\eta\). It is therefore unchanged as well.